% =====================================================================
%  CS 6330 - Introduction to Data Science
%  Homework 2: Data and Distributions
%  Author: Taek Lim
%
%  Build:  latexmk -pdf hw2_taeklim_report.tex
%      or: pdflatex hw2_taeklim_report.tex   (run twice for references)
%
%  Figure paths are relative to this file: ../figures/*.png
%  Figures and every number below come from ../code/p1.py and ../code/p2.py
%  (run from the repo root: uv run python hw2/code/p1.py).
% =====================================================================
\documentclass[11pt]{article}

\usepackage[letterpaper,margin=0.75in]{geometry}
\usepackage[T1]{fontenc}
\usepackage[utf8]{inputenc}
\usepackage{amsmath}
\usepackage{booktabs}     % \toprule \midrule \bottomrule
\usepackage{graphicx}
\usepackage{caption}
\usepackage{listings}
\usepackage{xcolor}
\usepackage[hidelinks]{hyperref}

\lstset{
  basicstyle=\ttfamily\footnotesize,
  breaklines=true,
  frame=single,
  rulecolor=\color{gray!40},
  backgroundcolor=\color{gray!5},
  showstringspaces=false,
  columns=fullflexible,
  literate={α}{{$\alpha$}}1 {λ}{{$\lambda$}}1 {μ}{{$\mu$}}1 {σ}{{$\sigma$}}1
           {≥}{{$\geq$}}1 {≤}{{$\leq$}}1,
}

\captionsetup{font=small, labelfont=bf}
\setlength{\parindent}{0pt}

% Layout follows the professor's answer template (00_doc/04_template_for_answers.pdf):
% bold headings, "-" items with the prompt text verbatim, answer under each item,
% one page each for Airport / Movie Rating / Problem 2.
\renewcommand{\labelitemi}{-}

% One figure inside an item: \answerfig{file}{caption}. \captionof (not a float) so it
% can sit inside itemize.
\newcommand{\answerfig}[2]{%
  \par\nopagebreak\smallskip % keep the prompt line on the same page as its figure
  {\centering
   \includegraphics[width=\linewidth]{../figures/#1.png}\par
   \captionof{figure}{#2}\par}
  \smallskip}

% Answer text under an item, set apart from the prompt line.
\newcommand{\answer}[1]{\par\smallskip{#1}\par\medskip}

\begin{document}

\hfill Taek Lim \quad CS 6330 Homework 2: Data and Distributions

\bigskip
% =====================================================================
\textbf{Problem 1.}

\medskip
\textbf{Airport data:}

\begin{itemize}
  \item What is the $\alpha$ if the distribution is power law?
  \answer{$\alpha = 1.612$, with $x_{\min} = 1$ route.
    Maximum-likelihood estimate $\hat\alpha = 1 + n / \sum_i \ln(x_i / x_{\min})$,
    $x_{\min} = \min_i x_i$, $n = 3409$.}

  \item What is the $\lambda$ if the distribution is exponential?
  \answer{$\lambda = 0.0504$. MLE $\hat\lambda = 1/\bar x = 1/19.845$.}

  \item What is $[a, b]$ if the distribution is uniform?
  \answer{$[a, b] = [1, 915]$, the minimum and maximum route counts (MLE).}

  \item What is $\mu$ and $\sigma$ if the distribution is normal?
  \answer{$\mu = 19.845$, $\sigma = 53.506$ (sample mean and sample standard deviation
    with $n-1$).}

  \item Screen shot of the visualization of distribution of data
  \answerfig{p1a_0_data}{Airport routes. Left: share of airports with exactly $k$ routes.
    Right: share with at least $x$ routes (CCDF). Both are on log-log axes. The median
    airport has 4 routes and 38\% have only 1 or 2, while the busiest has 915. The mean
    (19.8) is five times the median and the skewness is 6.0, so a few hubs hold most of
    the routes.}

  (These can be CDF, PDF, or QQplot)

  For each model below I drew $n = 3409$ random values from the fitted distribution
  (seed 6330). Left panel: CCDF of the data and of the model sample on log-log axes.
  Right panel: QQ plot of model-sample quantiles against data quantiles; points on the
  dashed line $y = x$ mean the model reproduces the data. $D$ is the two-sample
  Kolmogorov--Smirnov distance (largest gap between the two CDFs; 0 = identical).

  \item Screen shot of the visualization of distribution of data if it was power law
  \answerfig{p1a_1_powerlaw}{Power law ($\alpha = 1.612$, $x_{\min} = 1$). The CCDF
    follows the data from 1 to about 60 routes, and the QQ points lie on $y = x$ over
    that range. Above about 100 routes the model keeps going and produces airports with
    $10^4$ to $10^5$ routes, while the data bends down and stops at 915. $D = 0.209$.}

  \item Screen shot of the visualization of distribution of data if it was exponential
  \answerfig{p1a_2_exponential}{Exponential ($\lambda = 0.0504$). The model's sample
    goes below 1 route, which the data never does, and its tail dies out at about 150
    routes. The QQ points are too flat: the model gives small airports too many routes
    and hubs too few. $D = 0.379$.}

  \item Screen shot of the visualization of distribution of data if it was uniform
  \answerfig{p1a_3_uniform}{Uniform on $[1, 915]$. A typical sampled airport has about
    460 routes, but 99\% of real airports have fewer than 285. The model's CCDF stays
    near 1 until the right edge, and the QQ points lie far above $y = x$ for every
    quantile. $D = 0.857$, the worst of the four.}

  \item Screen shot of the visualization of distribution of data if it was normal
  \answerfig{p1a_4_normal}{Normal ($\mu = 19.8$, $\sigma = 53.5$). With $\sigma$ almost
    three times $\mu$, 36\% of the model's mass is below zero, which is impossible for a
    route count. The log axes can only show the positive part. The model's tail is also
    thin and ends near 200 routes. $D = 0.358$.}

  \item What is the distribution of data based on your observation? Discuss
  \answer{%
    \begin{center}
    \begin{tabular}{lcccc}
      \toprule
      & power law & exponential & uniform & normal \\
      \midrule
      KS distance $D$ (data vs model sample) & \textbf{0.209} & 0.379 & 0.857 & 0.358 \\
      \bottomrule
    \end{tabular}
    \end{center}

    \textbf{The airport data follows a power law, with a cutoff in the far tail.}
    It is the only one of the four models whose CCDF is a straight line on log-log
    axes. The data's CCDF is close to straight for more than two decades (1 to about
    100 routes), and the power-law sample matches it there. The other three models fail
    for structural reasons, not because a parameter is slightly off:
    \begin{itemize}
      \item \textbf{Uniform} says every route count from 1 to 915 is equally likely,
            but most airports are tiny.
      \item \textbf{Normal} is symmetric and allows negative counts, while the data is
            bounded below at 1 and skewed to the right (skewness 6.0).
      \item \textbf{Exponential} has the right shape (many small, few large), but its
            tail falls off too fast. With $\lambda = 0.05$, an airport with 500 routes
            has probability about $e^{-25}$, yet the data has five airports with more
            than 500.
    \end{itemize}
    The power law is not a perfect fit. $\hat\alpha = 1.61$ is below 2, so the fitted
    model has an infinite mean and its samples contain airports with $10^5$ routes. Real
    airports cannot grow that large: there are only 3409 airports to fly to, and
    runways, gates and slots limit the busiest hubs. That is why the data's CCDF bends
    down above about 100 routes. The best description is a power law with an upper
    cutoff, which is typical of networks that grow by preferential attachment: new
    routes go mostly to airports that already have many routes.}
\end{itemize}

% =====================================================================
\clearpage
\textbf{Movie Rating data:}

\begin{itemize}
  \item What is the $\alpha$ if the distribution is power law?
  \answer{$\alpha = 1.851$, with $x_{\min} = 1.9$ (the lowest average vote). Same MLE
    as for the airports, $n = 4392$.}

  \item What is the $\lambda$ if the distribution is exponential?
  \answer{$\lambda = 0.1606$. MLE $\hat\lambda = 1/\bar x = 1/6.227$.}

  \item What is $[a, b]$ if the distribution is uniform?
  \answer{$[a, b] = [1.9, 8.5]$, the minimum and maximum votes (MLE).}

  \item What is $\mu$ and $\sigma$ if the distribution is normal?
  \answer{$\mu = 6.227$, $\sigma = 0.893$ (sample mean and sample standard deviation
    with $n-1$).}

  \item Screen shot of the visualization of distribution of data
  \answerfig{p1b_0_data}{Movie votes. Left: histogram with one bin per possible vote
    (votes are rounded to 0.1). Right: empirical CDF. The distribution has one peak and
    is nearly symmetric around 6.3 (median 6.3, mean 6.23). The left tail is somewhat
    longer (skewness $-0.48$): a few movies are rated 2 to 4, and none are above 8.5.}

  (These can be CDF, PDF, or QQplot)

  For each model below I drew $n = 4392$ random values from the fitted distribution.
  Left panel: PDF (histogram) of the data and of the model sample. Right panel: QQ plot
  as for the airports. $D$ is the KS distance.

  \item Screen shot of the visualization of distribution of data if it was power law
  \answerfig{p1b_1_powerlaw}{Power law ($\alpha = 1.851$, $x_{\min} = 1.9$). The model
    puts its highest density at the lowest vote, which is where the data has almost
    nothing, and 13\% of its sample lies above 20, which is not a possible rating. The
    QQ plot is on log axes because the sample goes up to about 1000. $D = 0.485$.}

  \item Screen shot of the visualization of distribution of data if it was exponential
  \answerfig{p1b_2_exponential}{Exponential ($\lambda = 0.1606$). The model also peaks at
    zero and has a long right tail (4\% of its sample is above 20). In the QQ plot the
    points cross $y = x$ at a steep angle, so the model's spread is far larger than the
    data's. $D = 0.490$.}

  \item Screen shot of the visualization of distribution of data if it was uniform
  \answerfig{p1b_3_uniform}{Uniform on $[1.9, 8.5]$. The model is flat, but the data is
    concentrated between 5 and 7.5. The QQ plot is an S-shape: the model's low quantiles
    are too low and its high quantiles too high. $D = 0.394$.}

  \item Screen shot of the visualization of distribution of data if it was normal
  \answerfig{p1b_4_normal}{Normal ($\mu = 6.227$, $\sigma = 0.893$). The sample traces
    the same bell as the data, and the QQ points lie on $y = x$ from about 4.5 to 7.5.
    The only departures are at the ends: the data's lowest votes are lower than the
    model predicts (longer left tail), and the data stops at 8.5 while the model
    continues. $D = 0.055$, by far the smallest.}

  \item What is the distribution of data based on your observation? Discuss
  \answer{%
    \begin{center}
    \begin{tabular}{lcccc}
      \toprule
      & power law & exponential & uniform & normal \\
      \midrule
      KS distance $D$ (data vs model sample) & 0.485 & 0.490 & 0.394 & \textbf{0.055} \\
      \bottomrule
    \end{tabular}
    \end{center}

    \textbf{The movie votes follow a normal distribution, approximately.}
    The data has one peak near the middle of its range and is close to symmetric. The
    normal model predicts 68.3\% of movies within $\pm1\sigma$ and 95.4\% within
    $\pm2\sigma$; the data has 70.1\% and 95.6\%. In the QQ plot the points lie on the
    identity line through the bulk of the data, and the KS distance is about one seventh
    of the next-best model's.

    The other three models fail for structural reasons. The power law and the
    exponential are both monotonically decreasing: their most likely value is the
    smallest one, so they predict mostly terrible movies and a long tail of ratings
    above 10. The uniform model ignores the peak altogether.

    The fit is good but not exact, for two reasons. First, the data is bounded to
    $[1, 10]$ and slightly left-skewed: bad movies can drop further below the average
    than good movies can rise above it. Second, votes are rounded to one decimal, which
    gives the QQ plot its small steps. A normal shape is expected because each value is
    an \emph{average} of many individual votes. By the central limit theorem an average
    of many independent ratings tends toward a normal distribution, whatever the
    distribution of each single rating.}
\end{itemize}

% =====================================================================
\clearpage
\textbf{Problem 2.}

\begin{itemize}
  \item Screenshot of distribution/histogram of the age
  \answer{How the ages were computed. The birth dates have the form \texttt{dd-Mon-yy}
    (a few use a 4-digit year). Every 2-digit year is read as 19\textit{yy}: IMDB-WIKI
    was collected in 2015, so a year such as \texttt{20}--\texttt{24} cannot mean
    2020--2024, and the counts per year decrease smoothly from the 1920s down to
    \texttt{00}, with no break that would mark a switch to the 2000s. The 129 rows with
    year \texttt{0000} (a missing-value placeholder) were dropped. Per the professor's
    instruction (2026-09-23), age is now the year difference \textbf{2026 minus birth
    year} (month/day ignored), kept only if $1 \le$ age $\le 100$. Keeping that range
    leaves \textbf{453{,}905} people (mean 56.87, median 55, range 27 to 100). The
    minimum of 27 is a result of the 19\textit{yy} rule, since the latest possible birth
    year is 1999 ($2026-1999=27$). Of the rest, 6{,}684 rows had age $>100$ and were
    dropped; none had age $<1$.}
  \answerfig{p2_1_age}{Age in 2026 ($n = 453{,}905$). One broad peak in the 40s and 50s
    and a long right tail up to 100; most of the mass is between 36 and 71.}

  \item Screenshot of distribution/histogram of the first digit of age
  \answerfig{p2_2_first_digit}{First digit of age, with the uniform level $1/9$ (dashed)
    and Benford's law (orange). The shares are 0.2\%, 0.7\%, 7.9\%, 24.5\%, 30.0\%,
    18.9\%, 9.9\%, 5.4\% and 2.5\% for digits 1 to 9. The data is far from uniform and
    also far from Benford: it peaks at 5 and 4, where Benford predicts 1 to be the most
    common.}

  \item Screenshot of distribution/histogram of the last digit of age
  \answerfig{p2_3_last_digit}{Last digit of age, with the uniform level $1/10$ (dashed).
    Every digit has a share between 8.95\% and 11.10\%, within 1.10 percentage points of
    10\%.}

  \item Discussion (what are those distribution?) Why do you think it is like this?
  \answer{%
    \textbf{The last digit is uniform (approximately); the first digit is not.}

    \textbf{Last digit.} This is expected. The last digit of an age is its value modulo
    10. Every block of ten ages (20--29, 30--39, \dots) contains each last digit exactly
    once. The age distribution changes slowly compared with a step of one year, so within
    any block of ten the ten ages have roughly equal counts, and summing over the blocks
    gives each last digit about 10\%. The remaining deviations of up to $\pm1.1$ points
    come from year-to-year noise in how many actors were born in a given year (the jagged
    bars in the age histogram), not from any preference for particular digits.

    \medskip
    \textbf{First digit.} This is also expected, once the two boundary decades are
    accounted for separately. For ages 30 to 99, the first digit \emph{is} the decade of
    life, and grouping by decade reproduces the shape of the age distribution: it peaks
    at 5 (the 50s, 30.0\%) and 4 (the 40s, 24.5\%), matching the median of 55. Digit 2 is
    small (0.7\%) not because few actors are in their 20s, but because the $1 \le
    \text{age} \le 100$ rule only keeps ages 27 to 29 in that decade, three values out of
    ten, a boundary effect of the reference year rather than a property of the data.
    Digit 1 is smaller still (0.2\%): no age from 10 to 19 survives the rule (the
    youngest possible age is 27), so the entire digit-1 count comes from the 816 people
    with age exactly 100 ($816 / 453{,}905 = 0.18\%$, matching the observed share). For
    the first digit to be uniform, every decade would need the same number of actors,
    which working actors clearly are not, and the two extreme digits are further skewed
    by where the $[1, 100]$ window happens to cut.

    \medskip
    \textbf{It is not Benford's law either.} Benford's law (first digit $d$ with
    probability $\log_{10}(1 + 1/d)$, so 1 is the most common at 30.1\%) applies to data
    that spans several orders of magnitude, such as populations, prices or the airport
    route counts in Problem~1, where the log of the value is spread out roughly evenly.
    Age spans less than one order of magnitude (27 to 100) and is concentrated around
    55, so its first digit shows where that one hump sits, not Benford's pattern. In short,
    the last digit is uniform because it throws away where the value is and keeps only
    its position within a block of ten. The first digit keeps exactly that information,
    so its distribution follows the shape of the age distribution.}
\end{itemize}

% =====================================================================
\bigskip
\textbf{Tooling.} Python 3 / pandas / numpy / scipy / matplotlib; all estimates computed
programmatically, none typed by hand; analysis code, figure code and write-up drafted
with Claude Code. Code: \texttt{hw2/code/p1.py} and \texttt{hw2/code/p2.py}, listed below.

\clearpage
\textbf{Code}

\medskip
\textbf{Problem 1:} \texttt{hw2/code/p1.py}
\lstinputlisting[language=Python]{../code/p1.py}

\textbf{Problem 2:} \texttt{hw2/code/p2.py}
\lstinputlisting[language=Python]{../code/p2.py}

\end{document}
